\documentclass[german,10pt]{book}      \input{preambel}         \begin{document}  \setcounter{chapter}{0}

\chapter{Superposition}
% Kap x
\label{chap_Superposition}

\info{Thomas Filk}{5.10.2025}%

Die Superposition ist ein zentraler Begriff der Quantentheorie. Er bezeichnet
die Eigenschaft, dass zu je zwei realisierbaren Zust\"anden in der Quantentheorie
auch Linearkombinationen dieser beiden Zust\"ande wiederum realisierbaren
Zust\"anden entsprechen. Da Zust\"ande in der Quantentheorie oft durch einen
normierten Vektor dargestellt werden -- wobei auch eine (quadrat-)normierte Wellenfunktion als
Vektor (in einem Vektorraum von Funktionen) interpretiert werden kann -- vergleicht man 
die Superposition von Zust\"anden insbesondere in der Schule gerne mit der vektoriellen 
Summe von Vektoren, wie sie beispielsweise
aus der klassischen Mechanik f\"ur Geschwindigkeiten oder Kr\"afte bekannt ist. 
Dies ist nicht immer unproblematisch, da es sich streng genommen bei
der mathematischen Darstellung quantenmechanischer Zust\"ande nicht um
Vektoren sondern um Strahlen bzw.\ 1-dimensionale Unterr\"aume eines Vektorraums
handelt. F\"ur diese hat die Superposition eine etwas andere Bedeutung als f\"ur
Vektoren.

\section{Zust\"ande werden durch Strahlen repr\"asentiert}

Die Schr\"odinger-Gleichung ist eine lineare Differenzialgleichung f\"ur
Wellenfunktionen:
\begin{equation}
     {\rm i} \hbar \frac{\partial}{\partial t} \psi(x,t) = \left(- \frac{\hbar^2}{2m} \Delta + V(x) \right) \psi(x,t)  \, .
\end{equation}
In etwas abstrakterer Notation handelt es sich um eine Differentialgleichung f\"ur
Vektoren in einem Hilbert-Raum (also -- etwas vereinfacht -- einem Vektorraum mit einem Skalarprodukt),
\begin{equation}
     {\rm i} \hbar \frac{\partial}{\partial t} | \psi(t) \rangle =  H | \psi(t) \rangle  \, . 
\end{equation}
$\psi(x,t) = \langle x | \psi(t) \rangle$ ist die Wellenfunktion
im Ortsraum zu einem zeitabh\"angigen Zustandsvektor $|\psi(t)\rangle$ und $H$ ist der Hamilton-Operator
$H=P^2/2m + V(Q)$, wobei in der Ortsraumdarstellung der Operator $Q$ einfach der Multiplikation
mit $x$ und der Operator $P$ dem Operator $-{\rm i}\hbar\nabla$ entspricht.

Die Linearit\"at der Schr\"odinger-Gleichung bedeutet, dass f\"ur zwei L\"osungen
$\psi_1(x,t)$ und $\psi_2(x,t)$ der Gleichung auch eine beliebige Linearkombination
$\psi(x,t)=\alpha \psi_1(x,t) + \beta \psi_2(x,t)$ ($\alpha,\beta \in \mathbb{C}$) eine
L\"osung der Gleichung ist. Zus\"atzlich zu der Eigenschaft, eine L\"osung der 
Schr\"odinger-Gleichung zu sein, fordern wir in der Physik noch
eine Normierungsbedingung, meist in der Form
\begin{equation}
              \int_{\mathbb{R}^n} |\psi(x,t) |^2 {\rm d}^n x = 1  \, .
\end{equation}
Der Grund, dass \"uberhaupt eine solche Normierungsbedingung zus\"atzlich gefordert
werden kann, liegt darin, dass zwei Wellenfunktionen $\psi_1(x)$ und $\psi_2(x)$, die sich 
nur multiplikativ um eine komplexe Zahl unterscheiden -- f\"ur die also 
$\psi_2(x)=\alpha \psi_1(x)~ (\alpha \in \mathbb{C}/\{0\})$ gilt --, 
denselben physikalischen Zustand definieren. Die mathematische
Repr\"asentation eines physikalischen Zustands entsprich also dem Strahl (bzw.\
1-dimensionalen Vektorraum) der durch einen Vektor $|\psi\rangle$ bzw.\ eine
Wellenfunktion $\psi(x)$ gegeben ist und aus allen Vielfachen $\alpha |\psi\rangle$ bzw.\
$\alpha \psi(x)$ besteht. (Die Zeitabh\"angigkeit spielt f\"ur das Folgende keine Rolle, daher
ber\"ucksichtigen wir sie nicht weiter.)

Ein Quantenzustand l\"asst sich mathematisch auch als
Projektionsoperator auf einen 1-dimensionalen Vektorraum (d.h.\ die Spur des Projektionsoperators
ist 1) darstellen. Dabei handelt es sich um den Vektorraum, der durch den Vektor $|\psi\rangle$ bzw.\
die Wellenfunktion $\psi(x)$ aufgespannt wird. Wurde $|\psi\rangle$ normiert -- gilt also
$\langle \psi|\psi \rangle =1$ --, ist der Projektionsoperator auf diesen Strahl 
\begin{equation}
                P_\psi = | \psi \rangle \langle \psi |   \, .
\end{equation}
(Wird der Zustand durch eine normierte Wellenfunktion $\psi(x)$ dargestellt, ist der Projektionsoperator
der Integrationskern $P(x,y)= \psi(x) \psi(y)^*$, und seine Anwendung auf eine beliebige andere
Wellenfunktion $\phi(x)$ ist durch $\int P(x,y) \phi(y)\, {\rm d}^n y$ gegeben.)

Der Projektionsoperator repr\"asentiert direkt den Zustand, d.h., eine freie Phase im Zustandsvektor
bzw.\ der Wellenfunktion f\"allt heraus. Es besteht also eine \glqq 1\,zu\,1\grqq-Beziehung zwischen der
Menge aller Quantenzust\"ande, die in einem Hilbert-Raum dargestellt werden k\"onnen, 
und der Menge der Projektionsoperatoren
mit der Spur 1. Definiert ist ein solcher Projektionsoperator durch folgende Bedingungen:
\begin{equation}
     P^\dagger = P  \hspace{1cm}  P^2= P \hspace{1cm}  {\rm Spur}\,P = 1  \, .
\end{equation}
$P^\dagger = P$ bedeutet, dass $P$ selbstadjungiert ist und somit reelle Eigenwerte
und (zu verschiedenen Eigenwerten) orthogonale Eigenr\"aume besitzt. Geometrisch
bedeutet diese Bedingung, dass die Projektion eines Vektors auf den durch $P$ beschriebenen
Strahl orthogonal zu diesem Strahl erfolgt, daher spricht man auch schon mal von einem
orthogonalen Projektionsoperator. $P^2=P$ ist die definierende Eigenschaft eines Projektionsoperators:
Wurde eine Projektion durchgef\"uhrt, hat eine weitere Projektion keine Wirkung; der projizierte
Vektor befindet sich schon in dem Strahl. Spur\,$P=1$ bedeutet, dass es sich um einen
Projektionsoperator auf einen 1-dimensionalen Vektorraum handelt. Im Allgemeinen gibt die
Spur die Dimension des Unterraums an, auf den $P$ projiziert.

\section{Superposition von Strahlen}

Projektionsoperatoren bzw.\ die 1-dimensionalen Unterr\"aume eines Hilbert-Raums
bilden selbst keinen Vektorraum. Und da quantenmechanische Zust\"ande durch
die 1-dimensionalen Unterr\"aume eines Hilbert-Raums dargestellt werden, gilt
f\"ur sie ebenfalls, dass sie keinen Vektorraum bilden und dass man sie im Allgemeinen
nicht addieren kann. Die Summe von zwei Projektionsoperatoren ist zwar wieder ein
linearer Operator, aber im Allgemeinen kein Projektionsoperator. Spurerhaltende positive
Linearkombinationen von zwei Projektionsoperatoren, d.h.\ $\alpha P_1 + (1-\alpha)P_2$
f\"ur $0 < \alpha < 1$, sind Dichtematrizen, die zwar auch physikalische Zust\"ande
beschreiben, allerdings keine reinen Zust\"ande sondern sogenannte gemischte
Zust\"ande. Diese gemischten Zust\"anden sollte man auf keinen Fall mit den reinen Zust\"anden
verwechseln, die man aus der Superposition von zwei Zustandsvektoren erh\"alt (siehe
auch Abschnitt \ref{Superposition_Gemisch}). 

Das Superpositionsprinzip hat f\"ur quantenmechanische Zust\"ande eine etwas andere
Bedeutung als f\"ur Vektoren: Zwei linear unabh\"angige 1-dimensionale Vektorr\"aume
definieren eine sogenannte lineare H\"ulle dieser beiden Vektorr\"aume. Allgemein besteht
die lineare H\"ulle zu einer Menge von Vektoren aus allen Vektoren, die sich als
Linearkombination dieser Vektoren schreiben lassen. Sind nun zwei (linear unabh\"angige) 
1-dimensionale Vektorr\"aume gegeben, kann man s\"amtliche Linearkombinationen aus 
den Vektoren dieser beiden Vektorr\"aume betrachten und diese bilden offensichtlich einen
2-dimensionalen Vektorraum. Jeder 1-dimensionale Unterraum dieser 2-dimensionalen
H\"ulle von zwei Quantenzust\"anden l\"asst sich als Superposition der beiden Quantenzust\"ande
interpretieren. 

Im Folgenden seien die beiden Vektorr\"aume, die Quantenzust\"ande darstellen sollen und
von denen wir Linearkombinationen betrachten wollen,
orthogonal. Das bedeutet, jeder Vektor aus dem einen Vektorraum ist orthogonal zu jedem
Vektor aus dem anderen Vektorraum. Das ist zwar nicht notwendig und viele der folgenden
\"Uberlegungen lassen sich auch auf nicht-orthogonale Vektorr\"aume verallgemeinern,
doch in nahezu allen F\"allen, die eine Schulrelevanz haben, ist diese Situation gegeben.
Der Grund ist, dass zwei Eigenzust\"ande zu einer quantenmechanischen Observablen
(repr\"asentiert durch einen selbstadjungierten Operator) immer orthogonal sind, sofern
die zugeh\"origen Eigenwerte, also die Messwerte der Observablen zu diesen Zust\"anden,
verschieden sind. Wenn also Superpositionen von \glqq echten Alternativen\grqq\ betrachtet
werden, werden diese Alternativen immer durch orthogonale Zustandsvektoren dargestellt.  

Wenn zwei orthogonale, normierte Vektoren (als Repr\"asentanten eines Quantenzustands) $|\psi_1\rangle$ 
und $|\psi_2\rangle$ gegeben sind
und wir eine normierte Linearkombination dieser beiden Vektoren betrachten, 
\begin{equation}
         |\psi \rangle = \alpha |\psi_1 \rangle + \beta |\psi_2 \rangle  \hspace{1cm}
           | \alpha |^2 + | \beta |^2 = 1 \, ,
\end{equation}
dann sind die Wahrscheinlichkeiten, bei einem System, das im Zustand $|\psi\rangle$
pr\"apariert wurde, bei einer Messung den Zustand $|\psi_1 \rangle$ bzw.\ den Zustand
$|\psi_2\rangle$ zu finden, durch $|\alpha|^2 = |\langle \psi_1|\psi \rangle|^2$ 
bzw.\ $|\beta|^2 =  |\langle \psi_2|\psi \rangle|^2$ gegeben. Dies ist die Born'sche Regel. 
Diese Wahrscheinlichkeiten entsprechen dem Kosinusquadrat des Winkels zwischen den
beiden Vektoren und damit dem Kosinusquadrat des Winkels zwischen den zugeh\"origen
Strahlen. Das bedeutet, sie h\"angen nicht von einem Repr\"asentanten (einem Vektor) auf
diesen Strahlen ab. Somit sind diese Wahrscheinlichkeiten f\"ur die Strahlen wohl definiert. 
Insbesondere h\"angen diese Wahrscheinlichkeiten nicht von der relativen Phase
zwischen der Linearkombination von zwei Vektoren ab, d.h., die Wahrscheinlichkeiten
sind f\"ur die Linearkombination
\begin{equation}
          | \psi \rangle =  \alpha |\psi_1\rangle + \beta {\rm e}^{{\rm i}\delta} |\psi_2 \rangle
\end{equation}
unabh\"angig von dem Winkel $\delta$, obwohl es sich f\"ur verschiedene Werte von $\delta$
nicht nur um verschiedene Vektoren $|\psi\rangle$ handelt sondern auch um verschiedene
Strahlen und damit verschiedene Zust\"ande. 

\section{Die Polarisation von Photonen als Beispiel}

Die Polarisation von Licht bzw.\ von einzelnen Photonen ist ein sch\"ones Beispiel
f\"ur einen Zustandsraum, der den Strahlen in einem Vektorraum entspricht.
Insbesondere die linearen Polarisationen, bei denen es keine Phasenverschiebung
zwischen den beiden Komponenten des elektrischen Feldst\"arkevektors gibt, lassen
sich als Strahlen in einem 2-dimensionalen reellen Vektorraum darstellen. Wenn
wir einen normierten Vektor $\vec{A}$ entlang einer Polarisationsrichtung als 
Repr\"asentanten der Polarisation w\"ahlen, entsprechen $\vec{A}$ und $-\vec{A}$
derselben Polarisation. 

\begin{figure}[htb]
\begin{picture}(80,60)(0,0)
\thicklines
\put(40,30){\makebox(0,0){$\bullet$}}
\put(10,30){\line(1,0){60}}
\end{picture}
\begin{picture}(80,60)(0,0)
\thicklines
\put(40,30){\makebox(0,0){$\bullet$}}
\put(40,0){\line(0,1){60}}
\end{picture}
\begin{picture}(80,60)(0,0)
\thicklines
\put(40,30){\makebox(0,0){$\bullet$}}
\put(20,10){\line(1,1){40}}
\end{picture}
\begin{picture}(80,60)(0,0)
\thicklines
\put(40,30){\makebox(0,0){$\bullet$}}
\put(60,10){\line(-1,1){40}}
\end{picture}
\begin{picture}(80,60)(0,0)
\multiput(10,30)(6,0){11}{\line(1,0){4}}
\put(57,34){\makebox(0,0){${\scriptsize \alpha}$}}
\thicklines
\put(40,30){\makebox(0,0){$\bullet$}}
\put(12,16){\line(2,1){58}}
\end{picture}
\caption{\label{fig_Superposition_1}%
Die vier linearen Polarisationen $h$ (horizontal), $v$ (vertikal), $+$ (positiv diagonal) 
und $-$ (negativ diagonal). Au\ss erdem ist eine allgemeine Polarisationsrichtung, gekennzeichnet
durch ihren Winkel $\alpha$ zwischen der Polarisationsrichtung und der Horizontalen, wiedergegeben.}
\end{figure}

Abbildung \ref{fig_Superposition_1} zeigt f\"unf lineare Polarisationen: $h$ steht f\"ur horizontale
Polarisation, $v$ f\"ur vertikale Polarisation, $+$ f\"ur positiv diagonale Polarisation (d.h., die
Polarisationsrichtung verl\"auft unter einem Winkel von $+45^\circ$ relativ zur horizontalen Achse 
entlang einer Diagonalen mit positiver Steigung) und $-$ f\"ur eine negativ
diagonale Polarisation (die Polarisationsrichtung verl\"auft unter einem Winkel von $-45^\circ$
entlang einer Diagonalen mit negativer
Steigung). Eine allgemeine Polarisationsrichtung k\"onnen wir durch ihren Winkel $\alpha$ relativ zur
horizontalen Richtung kennzeichnen. In der Bra-Ket-Schreibweise dr\"ucken wir formal diese Polarisationen durch
normierte Vektoren aus:
\begin{equation}
  |h\rangle = \left(\!\!  \begin{array}{c} 1 \\ 0 \end{array} \!\! \right) \hspace{0.6cm}
  |v\rangle = \left( \!\! \begin{array}{c} 0 \\ 1 \end{array} \!\! \right) \hspace{0.6cm}
  |+\rangle = \frac{1}{\sqrt{2}} \left(\!\!  \begin{array}{c} 1 \\ 1 \end{array} \!\! \right) \hspace{0.6cm}
  |-\rangle =  \frac{1}{\sqrt{2}} \left( \!\! \begin{array}{c} 1 \\ -1 \end{array} \!\! \right) \hspace{0.6cm}
  |\alpha\rangle = \left( \!\! \begin{array}{c} \cos \alpha \\ \sin \alpha \end{array} \!\! \right)
\end{equation}
Diese Vektoren sind jedoch nur Repr\"asentanten der Strahlen. Au\ss erdem wurden in diesem Fall
die Polarisationen $h$ und $v$ ausgezeichnet, indem $|h\rangle$ und $|v\rangle$ als die Basisvektoren
gew\"ahlt wurden.\footnote{Einfache Buchstaben oder Symbole wie $h$, $v$ oder $\psi$ beschreiben
allgemein einen quantenmechanischen Zustand; ein Ket-Vektor $|h\rangle, |v\rangle, |\psi\rangle$
bezeichnet einen normierten Vektor, der den Zustand im Hilbert-Raum repr\"asentiert.}

F\"ur diese Vektoren gilt:
\begin{equation}
        |+\rangle = \frac{1}{\sqrt{2}} \big( |h\rangle + |v \rangle \big)  \hspace{1cm} {\rm und} \hspace{1cm}
        |-\rangle = \frac{1}{\sqrt{2}} \big( |h\rangle - |v \rangle \big) 
\end{equation}
Umgekehrt gilt aber auch:
\begin{equation}
        |h\rangle = \frac{1}{\sqrt{2}} \big( |+\rangle + |- \rangle \big)  \hspace{1cm} {\rm und} \hspace{1cm}
        |v\rangle = \frac{1}{\sqrt{2}} \big( |+\rangle - |- \rangle \big) 
\end{equation}
F\"ur eine allgemeine Polarisation erhalten wir:
\begin{equation}
        |\alpha \rangle = \big( \cos \alpha |h\rangle + \sin \alpha |v \rangle \big) =
           \frac{1}{\sqrt{2}} \Big( \big( \cos \alpha + \sin \alpha \big) |+\rangle + 
                  \big( \cos \alpha - \sin \alpha \big) |- \rangle \Big) \, .
\end{equation}
Ein Vorteil der Bra-Ket-Schreibweise ist, dass sie unabh\"angig von der Wahl eines Koordinatensystems
ist. Die obigen Beziehungen zwischen verschiedenen Polarisationsrichtungen gelten in allen 
Koordinatensystemen, wohingegen die Darstellung durch Spaltenvektoren die Wahl einer Basis 
voraussetzt. 

Die Born'sche Regel wird bei Polarisationszust\"anden zu dem Gesetz von Malus:
\begin{equation}
          |  \langle \beta | \alpha \rangle |^2 = \cos^2 (\beta-\alpha) \, .
\end{equation}

\section{Superposition und Gemisch}
\label{Superposition_Gemisch}

Zwei wichtige Konzepte, die leider oft verwechselt werden, beziehen sich auf den
Unterschied zwischen Superposition und Gemisch. Sind zwei reine Zust\"ande
gegeben, so ist eine Superposition wieder ein reiner Zustand. Wenn die reinen (orthogonalen)
Zust\"ande $\psi_1$ und $\psi_2$ durch die normierten Vektoren $|\psi_1\rangle$
und $|\psi_2\rangle$ repr\"asentiert werden, dann repr\"asentiert auch der
Vektor
\begin{equation}
       | \psi \rangle = a |\psi_1\rangle + b| \psi_2\rangle  \hspace{1cm}
              {\rm mit} ~~  |a|^2 + |b|^2 = 1  
\end{equation}
wieder einen reinen Zustand. Die Koeffizienten $a$ und $b$ k\"onnen dabei bis auf
ihre Normierung beliebig sein und kennzeichnen verschiedene Superpositionen
von $|\psi_1\rangle$ und $|\psi_2\rangle$. Messungen, die zwischen $\psi_1$ und
$\psi_2$ unterscheiden, liefern mit einer Wahrscheinlichkeit (oder relativen H\"aufigkeit) 
$|a|^2$ den Zustand $\psi_1$ und mit der Wahrscheinlichkeit (bzw.\ relativen H\"aufigkeit) $|b|^2$
den Zustand $\psi_2$. Findet man jedoch eine Messvorschrift, sodass die
Superposition $| \psi \rangle = a |\psi_1\rangle + b| \psi_2\rangle$ ein Eigenzustand des
zugeh\"origen selbstadjungierten Operators ist, dann findet man dieses Ergebnis
auch bei jeder einzelnen Messung.

Man kann die Zust\"ande $\psi_1$ und $\psi_2$ auch durch ihre Projektionsoperatoren
darstellen: $P_1=|\psi_1\rangle \langle \psi_1|$ und $P_2=|\psi_2\rangle \langle \psi_2|$.
Bildet man nun eine Linearkombination der Form
\begin{equation}
         \rho = p_1 P_1 + p_2 P_2  \hspace{1cm} {\rm mit} ~~ p_1+p_2=1 ~~(p_i > 0) \, ,
\end{equation}
so beschreibt $\rho$ ein Gemisch aus den beiden Zust\"anden $\psi_1$ und $\psi_2$. 
In einem Ensemble von physikalischen Systemen, das dieses Gemisch repr\"asentiert, sind Systeme 
im Zustand $\psi_1$
mit einer relativen H\"aufigkeit $p_1$ und Systeme im Zustand $\psi_2$ mit der
relativen H\"aufigkeit $p_2$ vertreten. Messungen, die zwischen $\psi_1$ und
$\psi_2$ unterscheiden, liefern mit einer Wahrscheinlichkeit (oder relativen H\"aufigkeit) 
$p_1$ den Zustand $\psi_1$ und mit der Wahrscheinlichkeit (bzw.\ relativen H\"aufigkeit) $p_2$
den Zustand $\psi_2$. Solche Messungen, k\"onnen also nicht zwischen einem Gemisch
und einer Superposition unterscheiden. In diesem Fall gibt es jedoch keine Messung,
bei der immer dasselbe Ergebnis herauskommt. $\rho$ ist kein Objekt, das man als
Eigenzustand zu einem selbstadjungierten Operator interpretieren kann. Das Fehlen eines
solchen Operators bzw.\ das Fehlen einer entsprechenden Messvorschrift ist das
charakteristische Zeichen f\"ur ein Gemisch. 

Betrachten wir hierzu nochmals das Beispiel der Polarisation von Photonen. Eine
Superposition eines Zustands, der ein Photon mit der Polarisation $|h\rangle$
beschreibt, mit einem Zustand, der ein Photon mit der Polarisation $|v\rangle$
beschreibt, ergibt
\begin{equation}
            |+\rangle = \frac{1}{\sqrt{2}}\big( |h\rangle + |v\rangle  \big) 
\end{equation}
(hier wurde die spezielle Superposition zu $a=b=1/\sqrt{2}$ gew\"ahlt). Dieser Vektor
beschreibt den Zustand eines Photons mit der Polarisation $+$, also eine unter $+45^\circ$
geneigte Polarisation. In einem Mach-Zehnder-Interferometer mit polarisierenden
Strahlteilern kann man diese Gleichung experimentell realisieren: Ein Photon im
Zustand $|+\rangle$ wird am ersten Strahlteiler, der die Achsen $|h\rangle$ und $|v\rangle$
haben soll, in die beiden Bestandteile $|h\rangle$ und $|v\rangle$ entsprechend der
obigen Formel aufgespalten. Am zweiten Strahlteiler k\"onnen diese beiden Anteile
bei exakt gleicher optischer Wellenl\"ange der beiden Wege innerhalb des 
Mach-Zehnder-Interferometers wieder in Superposition gebracht werden und man
erh\"alt hinter dem zweiten Strahlteiler wieder ein diagonal polarisiertes Photon.

Man kann zur Demonstration auch die koh\"arente \"Uberlagerung von zwei 
Laserstrahlen verwenden. Wenn die Phasen der beiden Laser richtig abgestimmt sind,
kann man die Strahlen mit einem geeigneten polarisationsabh\"angigen Strahlteiler
so \"uberlagern, dass man einen diagonal polarisierten Strahl (der doppelten Intensit\"at)
erh\"alt. 

Misst man die Photonen (oder den Laserstrahl) mit einem Polarisationsstrahlteiler
zu den Achsen $h$ und $v$, erh\"alt man bei den Photonen in der H\"alfte der 
F\"alle $h$ und in der H\"alfte der F\"alle $v$, beim Laserstrahl erh\"alt man in beide
Richtungen jeweils die halbe Intensit\"at. Stellt man jedoch den Polarisationsstrahlteiler
auf $+$- bzw.\ $-$-diagonal, findet man alle Photonen (bzw.\ die gesamte Laserintensit\"at)
in der $+$-Richtung. Dies ist das Zeichen, dass eine reine Polarisation $+$ vorliegt.

Ein gemischter Zustand von Photonen, die zur H\"alfte die Polarisation $h$ und
zur H\"alfte die Polarisation $v$ haben, wird bei einem $h/v$-orientierten Strahlteiler
ebenfalls zur H\"alfte das Ergebnis $h$ und zur H\"alfte das Ergebnis $v$ liefern.
Dasselbe gilt aber auch, wenn man den Strahlteiler in eine beliebige andere
Richtung orientiert. Es gibt keine Position, bei der alle Photonen in dieselbe
Richtung abgelenkt werden. Dies ist wieder das Signal f\"ur einen gemischten
Zustand. (Streng genommen m\"usste man bei diesem Test auch auf zirkulare
Polarisation testen, die bei linearen Polarisationsstrahlteilern ebenfalls f\"ur
alle Richtungen ein 50\% -50\%--Ergebnis liefern w\"urde.) 

\section{Vergleich: Superposition und Vektoraddition}

Offenbar haben Superposition und Vektoraddition vieles gemein. Insbesondere l\"asst sich die
Superposition als \glqq Vektoraddition von Repr\"asentanten\grqq\ darstellen. F\"ur einen besseren
Vergleich betrachten wir daher nochmals das Paradebeispiel aus der Physik, bei dem die
Vektoraddition Anwendung findet: das Kr\"afteparallelogramm.

Zu zwei Kr\"aften $\vec{F}_1$ und $\vec{F}_2$, die am selben Punkt eines K\"orpers angreifen,
k\"onnen wir eine resultierende Kraft $\vec{F}=\vec{F}_1 + \vec{F}_2$ definieren, die sich als
die vektorielle Summe der beiden Kr\"afte darstellen l\"asst (in der Schule spricht man 
h\"aufig vom \glqq Parallelogrammgesetz\grqq). Die physikalische Wirkung, die diese Kraft
$\vec{F}$ auf den K\"orper aus\"ubt, ist dieselbe, wie die physikalische Wirkung der beiden
Kr\"afte $\vec{F}_1$ und $\vec{F}_2$. Aus der Beobachtung der Bewegung des K\"orpers
alleine lassen sich die beiden F\"alle nicht unterscheiden.

Umgekehrt k\"onnen wir eine gegebene Kraft $\vec{F}$ immer nach zwei vorgegebenen
(linear unabh\"angigen) Richtungen zerlegen und erhalten auf diese Weise zwei
Kr\"afte $\vec{F}_1$ und $\vec{F}_2$, die entlang der vorgegebenen Richtungen zeigen
und deren vektorielle Summe gleich der Kraft $\vec{F}$ ist.\footnote{Der Einfachheit halber
betrachten wir hier Kr\"afte in einer Ebene. Ansonsten m\"ussen die beiden Richtungen so
gew\"ahlt werden, dass die Kraft $\vec{F}$ in der von ihnen aufgespannten Ebene liegen.}
Da diese beiden Richtungen
(bis auf die lineare Unabh\"angigkeit) beliebig gew\"ahlt werden k\"onnen, l\"asst sich jede
beliebige Kraft auf unendlich viele verschiedenen Weisen als Summe von zwei anderen
Kr\"aften schreiben. 

Vieles davon l\"asst sich auf die Superposition von Zust\"anden \"ubertragen: Jeder
beliebige Zustand l\"asst sich auf unendlich viele verschiedene Weisen als 
Superposition von anderen Zust\"anden schreiben. Allerdings lassen sich zu je zwei
Zust\"anden auch unendlich viele Zust\"ande als Superposition dieser beiden Zust\"ande
schreiben. Eindeutig wird die Superposition erst, wenn wir zu den beiden Zust\"anden
Einheitsvektoren als Repr\"asentanten w\"ahlen und die Koeffizienten der Linearkombination
festlegen. 

Andererseits gibt es aber auch wichtige Unterschiede, die gerade bei Sch\"ulerinnen und
Sch\"ulern leicht zu Fehlvorstellungen f\"uhren k\"onnen. So gibt es beispielsweise zu
jedem Vektor den entsprechenden negativen Vektor, und die vektorielle Summe dieser
beiden Vektoren ergibt den Nullvektor. Dies spielt eine wichtige Rollen, wenn wir
Systeme im Kr\"aftegleichgewicht betrachten. W\"ahlen wir jedoch einen Vektor
als Repr\"asentanten eines Zustands, dann ist der dazu negative Vektor ebenfalls ein
Repr\"asentant desselben Zustands. In diesem Fall ist es sinnlos, eine Superposition
dieser beiden Vektoren zu betrachten. 

Ein weiterer wichtiger Unterschied zwischen Superposition und vektorieller Summe
besteht in der Normierung. Wenn wir zwei Vektoren, beispielsweise Kr\"afte, addieren,
beschreibt die Summe einen Vektor, dessen L\"ange sich nach dem Parallelogrammgesetz
bestimmt. Diese L\"ange ist wesentlich, ansonsten beschreibt die vektorielle Summe
keine Kraft, welche die urspr\"unglichen beiden Kr\"afte ersetzen kann. Bei der
Repr\"asentation der Superposition von Zust\"anden als vektorielle Summe muss jedoch
der resultierende Vektor wieder auf die L\"ange 1 normiert werden. Ansonsten beschreibt
er (im \"ublichen Sinne) keinen Zustand. Etwas anders ausgedr\"uckt: W\"ahrend bei
Vektoren, z.B.\ bei Kr\"aften, die L\"ange der Vektoren wesentlich ist, ist beim Zustand
(und der Superposition von Zust\"anden) eigentlich nur der Winkel zwischen den
Vektoren relevant. Vektoren, die sich in ihrer L\"ange unterscheiden, beschreiben
denselben Zustand.

\section{H\"aufige Fehlvorstellungen zur Superposition}

Man h\"ort sehr oft, dass ein bestimmter Zustand ein \glqq Superpositionszustand\grqq\ sei.
Das klingt so, als ob es auch Zust\"ande g\"abe, die keine Superpositionszust\"ande sind.
Hinter dieser Sprechweise steckt meist die implizite Annahme, dass gewisse Zust\"ande
ausgezeichnet sind. Diese Sprechweise wird beispielsweise oft verwendet, wenn man 
\"uber die Superposition von Zust\"anden mit klassischen Eigenschaften spricht.
Beispielsweise im Zusammenhang mit Schr\"odingers Katze bezeichnet man die
Zust\"ande $|{\rm tot}\rangle$ und $|{\rm lebendig}\rangle$ als klassisch realisierbare
Zust\"ande, wohingegen eine Superposition von $|{\rm tot}\rangle$ und $|{\rm lebendig}\rangle$
als Superpositionszustand bezeichnet wird. Ebenso spricht man im Zusammenhang mit
dem Messproblem gerne von einem Superpositionszustand, wenn man die Superposition
von zwei Zeigerstellungen des Messinstruments betrachtet (die typischerweise nie
beobachtet wird). 

Ein anderes Beispiel, bei dem gerne von Superpositionszust\"anden gesprochen wird, ist
die Quanteninformation. Die Zust\"ande $|0\rangle$ und $|1\rangle$ sind ausgezeichnet,
weil sie in der Quanteninformation die typischen Basiszust\"ande sind und weil
abschlie\ss ende Messungen an einem Quantenzustand sich auf diese Basis beziehen.
Superpositionszust\"ande sind dann Zust\"ande, die nicht diesen beiden ausgezeichneten
Basiszust\"anden entsprechen. 

Ohne die Angabe einer ausgezeichneten Basis (oder einer ausgezeichneten Observablen,
die eine solche Basis definiert) ist es sinnlos, von einem Superpositionszustand zu
sprechen oder gar einen Zustand als \glqq keine Superposition\grqq\ zu bezeichnen. In
diesem Zusammenhang ist es sogar ein Fehler, von reinen und gemischten
Zust\"anden zu sprechen: Auch die Superpositionszust\"ande zu einer ausgezeichneten
Basis sind reine Zust\"ande.

Eine weitere verbreitete Fehlvorstellung liegt in der Verwechslung von Superposition
und Gemisch. Hierbei handelt es sich um zwei vollkommen verschiedene Konzepte:
Eine Superposition von zwei reinen Zust\"anden ist wieder ein reiner Zustand und
entspricht einem Strahl in einem Hilbert-Raum. Ein Gemisch von zwei Zust\"anden
ist kein reiner Zustand und l\"asst sich nicht als Strahl in einem Hilbert-Raum darstellen,
sondern wird meist durch eine Dichtematix dargestellt. In diesem Fall liegen physikalisch
(mindestens) zwei verschiedene Zust\"ande vor und die Eigenwerte der Dichtematrix
geben an, mit welcher Wahrscheinlichkeit diese Zust\"ande bei einer willk\"urlichen
Wahl aus einem Ensemble physikalischer Systeme, das diesem Gemisch entspricht, gezogen werden.

\end{document}

